% Korean edition of figs/sectors.tex.
\fig{sectors}{섹터 한도에 가장 가까운 것은 경기소비재로, 2.01\%p 남았다}[섹터 비중(\%)과 한도, 여유(\%p), 보유 비중이 큰 순서. 2026년 11월 24일 09:52 KST]{%
{\tabfont\begin{tabularx}{\linewidth}{>{\RaggedRight\arraybackslash}X R{19mm} R{15mm} R{14mm} R{16mm}}
\thd{GICS 섹터} & \thd{시장 비중} & \thd{포트폴리오} & \thd{한도} & \thd{여유}\\
45 IT & 52.90 & 45.74 & 105.80 & 60.06\\\rowsep
20 산업재 & 17.70 & 30.17 & 35.40 & 5.23\\\rowsep
\hl 25 경기소비재 & 5.75 & 9.49 & 11.50 & 2.01\\\rowsep
50 커뮤니케이션 서비스 & 5.00 이하 & 6.79 & 10.00 & 3.21\\\rowsep
35 헬스케어 & 5.69 & 6.54 & 11.38 & 4.84\\
\end{tabularx}}}{음영 행: 여유가 가장 적은 섹터. 한도 = max(2 $\times$ 시장 비중, 10\%). 여유 = 한도 $-$ 포트폴리오 비중. 가장 큰 섹터인 IT는 60.06\%p가 남았다. 시장 비중은 한도에서 역산했다. 커뮤니케이션 서비스는 10\% 하한이 적용되므로 시장 비중의 상한만 알 수 있다. 이 5개 섹터가 포트폴리오의 98.73\%를 차지한다.}{\code{guidelines()} 출력, DART Quant Team}
